\section{Introduction}
\label{sec:intro}

Long-lived edge data needs confidentiality and integrity without assuming a datacenter-class accelerator.  Post-quantum key establishment is relevant to that lifetime: ML-KEM establishes a shared secret, while symmetric authenticated encryption remains the appropriate mechanism for bulk data~\cite{mlkem}.  SHA-256 can bind integrity-tree material to the persisted format~\cite{sha256}.  In practice, however, a secure-storage prototype must first get the ordinary systems details right: the key hierarchy must survive remount, ciphertext must become durable before it is published by metadata, and every claimed accelerator path must have a reproducible measurement.  Existing filesystem encryption is already mature~\cite{fscrypt}, as are user-space encrypted filesystems~\cite{gocryptfs}; a new user-space design must therefore earn complexity rather than merely restate it.  A trusted execution environment such as OP-TEE supplies a different protection boundary~\cite{optee}, not a substitute for file-format recovery logic.

This paper presents AEGIS-Q, an architecture that redefines edge storage not through theoretical cryptography, but through rigorous systems engineering. While previous work has struggled to balance secure storage with heavy AI inference loads on Unified Memory Architectures (UMA) like the NVIDIA Jetson, AEGIS-Q explores a placement policy that is evaluated only on the retained artifact set. The current revision does not claim guaranteed Quality of Service (QoS) or a verified zero-copy Direct Memory Access (DMA) path.

We intentionally scope out theoretical physical security claims, such as oscilloscope-based side-channel resistance, recognizing that such metrics distract from the core systems challenge: avoiding resource starvation in a multitenant edge environment. Instead, AEGIS-Q focuses on solving real bottleneck problems in the edge data path.

The paper makes four core systems contributions:
\begin{enumerate}
\item \textbf{Placement-Scoped GPU Path:} The implementation keeps GPU work inside managed memory with explicit prefetch and stream synchronization. It does not claim verified \texttt{O\_DIRECT} NVMe-to-GPU DMA or a zero-copy storage pipeline.
\item \textbf{No Verified Fast-Notification Path:} The revised paper does not claim a completed \texttt{eBPF} or \texttt{io\_uring} completion-bypass mechanism, and it does not treat those techniques as evaluated contributions.
\item \textbf{Telemetry Not Yet Closed:} The repository contains a prototype telemetry-aware scheduler and measured TensorRT/YOLO interference artifacts, but this revision does not claim a validated \texttt{tegrastats}-driven QoS controller.
\item \textbf{Transparent Scope Boundaries:} The revision explicitly distinguishes verified prototype behavior from the retained matched \texttt{fscrypt} and \texttt{dm-crypt} baselines, while leaving broader baseline expansion as future work.
\end{enumerate}

By abandoning unverifiable security claims and focusing on the verified prototype behavior, AEGIS-Q establishes an evidence-backed foundation for future edge-storage evaluations.  Section~\ref{sec:background} grounds the storage and UMA assumptions, Section~\ref{sec:related_work} positions the design against nearby storage and accelerator systems, Section~\ref{sec:design} describes the implemented planes and trust boundaries, Section~\ref{sec:design_impl} records the implementation choices that are easy to misread from a diagram, Section~\ref{sec:evaluation} reports the measured primitive placement and negative controls, and the remaining sections discuss limitations, conclusions, and reproducibility artifacts.
